\documentclass[11pt]{article}

\usepackage[margin=1in]{geometry}
\usepackage{booktabs}
\usepackage{array}
\usepackage{hyperref}
\usepackage{xcolor}

\title{Weekly Progress Report: Reproducing LA-CDM with Llama on MIMIC-IV}
\author{Omer Kahveci}
\date{September 2026}

\begin{document}

\maketitle

\section{Project Objective}

The objective of this project is to reproduce the Language Agent Clinical
Decision-Making (LA-CDM) pipeline using MIMIC-IV clinical data. The original
pipeline was adapted to use Meta-Llama-3-8B-Instruct instead of Qwen because
the required Chinese-developed models are not available for use on HiPerGator.
Consequently, this work is a methodological reproduction of LA-CDM rather than
an exact reproduction of the paper's model-specific results.

\section{Cohort Construction}

The source data were MIMIC-IV version 2.2 and MIMIC-IV-Note version 2.2. A
candidate cohort was extracted for four acute abdominal conditions:
appendicitis, cholecystitis, diverticulitis, and pancreatitis. The extracted
tables were then processed using the official MIMIC Clinical Decision-Making
Dataset construction code.

The candidate extraction contained 12,441 admissions from 9,336 unique
patients. The official cohort builder reduced this set to 2,400 eligible
admissions. Table~\ref{tab:cohort} reports the final validated cohort.

\begin{table}[ht]
    \centering
    \caption{Final cohort composition.}
    \label{tab:cohort}
    \begin{tabular}{lr}
        \toprule
        Condition & Cases \\
        \midrule
        Appendicitis & 957 \\
        Cholecystitis & 648 \\
        Diverticulitis & 257 \\
        Pancreatitis & 538 \\
        \midrule
        \textbf{Total} & \textbf{2,400} \\
        \bottomrule
    \end{tabular}
\end{table}

Validation confirmed that the cohort contained no invalid records, diagnosis
leakage, or overlap in admission identifiers between disease groups. Every
case contained a patient history, physical examination, laboratory data, and
abdominal imaging. The laboratory mapping contained 1,104 rows.

\section{Dataset Preparation}

The cohort was divided into patient-level training, validation, and test sets,
as shown in Table~\ref{tab:splits}. Meta-Llama-3-8B-Instruct was then used to
create deterministic summaries of all 2,400 patient histories. Summary
generation was made resumable and checkpointed after every 25 newly completed
rows.

\begin{table}[ht]
    \centering
    \caption{Dataset split sizes.}
    \label{tab:splits}
    \begin{tabular}{lr}
        \toprule
        Split & Cases \\
        \midrule
        Training & 1,920 \\
        Validation & 240 \\
        Test & 240 \\
        \midrule
        \textbf{Total} & \textbf{2,400} \\
        \bottomrule
    \end{tabular}
\end{table}

The prepared data passed validation with no blank summaries, duplicate
admission identifiers, malformed laboratory fields, malformed radiology
fields, invalid labels, or overlap between splits.

\section{Computing and Model Configuration}

Experiments were conducted on HiPerGator using two NVIDIA L4 GPUs and 64~GB of
system memory. One GPU hosts the trainable Transformers model, while the second
hosts the vLLM inference engine. The model is Meta-Llama-3-8B-Instruct in local
Transformers format. Training uses bfloat16 and LoRA adapters with rank 8,
scaling parameter 16, and dropout 0.1.

The production GRPO configuration generates eight trajectories per example.
The eight generations are divided into sequential microbatches for the forward
and backward passes. Production training uses a microbatch size of one after a
longer trajectory exceeded L4 memory with a microbatch size of two. This preserves the required GRPO comparison
group while remaining within the 22~GB memory capacity of each L4 GPU.

\section{Engineering and Reproducibility Work}

Several changes were required to run the original pipeline reliably with
Llama and the HiPerGator software environment:

\begin{itemize}
    \item Created compatible dependency pins for PyTorch, vLLM, Transformers,
          XGrammar, Numba, and LLVMlite.
    \item Moved the Conda environment and package cache to Blue storage after
          the home filesystem quota was exhausted.
    \item Added ordered extraction, construction, and validation scripts.
    \item Added Slurm jobs with job-specific logs and output directories.
    \item Corrected duplicate aliases and missing synonym rows in the generated
          laboratory-test mapping.
    \item Configured the Llama EOS token for padding.
    \item Removed an unnecessary reference-model copy during evaluation, which
          eliminated a full model's worth of GPU memory use.
    \item Updated the action parser to tolerate cosmetic Llama formatting while
          continuing to validate actions against strict test and disease lists.
    \item Corrected the evaluation pipeline so custom accuracy, calibration,
          reward, test-count, and diagnostic-cost metrics are saved.
    \item Made evaluation repeatable by supplying a fixed vLLM seed and using
          near-greedy decoding with temperature 0.01.
\end{itemize}

\section{Zero-Shot Baseline}

The full zero-shot evaluation covered all 240 held-out test cases and completed
in 2,540.7 seconds, or approximately 42.3 minutes. The results are summarized
in Table~\ref{tab:baseline}.

\begin{table}[ht]
    \centering
    \caption{Meta-Llama-3-8B-Instruct zero-shot baseline results.}
    \label{tab:baseline}
    \begin{tabular}{lr}
        \toprule
        Metric & Result \\
        \midrule
        Overall diagnostic accuracy & 8.75\% \\
        Undiagnosed or invalid fraction & 87.92\% \\
        Accuracy among completed diagnoses & 72.41\% \\
        Macro F1, all cases & 15.54\% \\
        Macro F1, completed diagnoses only & 68.50\% \\
        Hypothesis-agent accuracy & 65.25\% \\
        Expected calibration error & 10.87\% \\
        Average tests requested & 1.16 \\
        Average diagnostic cost & \$798.08 \\
        Evaluation runtime & 42.3 minutes \\
        \bottomrule
    \end{tabular}
\end{table}

Table~\ref{tab:disease} reports disease-specific zero-shot performance when
undiagnosed cases are counted as incorrect.

\begin{table}[ht]
    \centering
    \caption{Zero-shot performance by diagnosis.}
    \label{tab:disease}
    \begin{tabular}{lrrrr}
        \toprule
        Condition & Support & Precision & Recall & F1 \\
        \midrule
        Appendicitis & 96 & 1.000 & 0.104 & 0.189 \\
        Cholecystitis & 65 & 0.800 & 0.062 & 0.114 \\
        Diverticulitis & 25 & 0.429 & 0.120 & 0.188 \\
        Pancreatitis & 54 & 0.571 & 0.074 & 0.131 \\
        \bottomrule
    \end{tabular}
\end{table}

The main zero-shot limitation was protocol completion rather than diagnostic
discrimination. Llama frequently requested an appropriate initial test and
produced clinically relevant reasoning, but then ended its response before
emitting the required \texttt{Action} and \texttt{Action Input} fields. Among
the 29 cases for which a valid diagnosis was completed, accuracy was 72.41\%.
However, the high undiagnosed fraction reduced overall accuracy to 8.75\%.

\section{Training Validation}

The training pipeline was tested incrementally before requesting a full
allocation. A two-step LoRA smoke test completed successfully with finite
losses and gradients, and KL divergence became nonzero after the first update.
The resulting adapter was saved and successfully reloaded for inference through
vLLM.

An initial production-scale test produced an out-of-memory error because its
temporary job configuration incorrectly set the forward microbatch size to
eight. After restoring the intended microbatch size of two, a one-step test
with all eight GRPO generations completed successfully. This verified the
production memory configuration.

An attempted two-epoch run revealed that the custom GRPO sampler expands this
configuration to 3,840 optimization steps, with an estimated runtime of about
75 hours. It also exceeded L4 memory on a longer trajectory at step five. The
revised training job therefore runs one complete 300-step schedule cycle: 100
Decision Agent GRPO steps, 100 Hypothesis Agent supervised fine-tuning steps,
and 100 Hypothesis Agent confidence-calibration GRPO steps. It uses
microbatches of one. A later long trajectory exposed a second memory peak in
the reference-log-probability calculation, which originally processed all
eight trajectories simultaneously. That no-gradient calculation is now also
chunked one trajectory at a time. Checkpoints are written every 25 steps.

\section{Problems and Resolutions}

\begin{table}[ht]
    \centering
    \caption{Major implementation problems and resolutions.}
    \label{tab:problems}
    \begin{tabular}{>{\raggedright\arraybackslash}p{0.35\textwidth}
                    >{\raggedright\arraybackslash}p{0.55\textwidth}}
        \toprule
        Problem & Resolution \\
        \midrule
        Conflicting package requirements & Constructed a mutually compatible,
        pinned environment instead of installing the original requirements
        file unchanged. \\
        Home filesystem quota exhausted & Stored the environment and package
        caches on Blue storage. \\
        PyTorch incompatible with Blackwell GPU architecture & Requested L4
        GPUs supported by the installed PyTorch build. \\
        Evaluation GPU out-of-memory & Removed the unnecessary reference-model
        copy and placed Transformers and vLLM on separate GPUs. \\
        Missing Llama padding token & Assigned the EOS token as the padding
        token and propagated its identifier to the model configuration. \\
        Qwen-specific output parsing & Added controlled tolerance for Markdown,
        whitespace, and prose reasoning while retaining action allowlists. \\
        Missing diagnostic metrics & Updated the Trainer-owned metrics mapping
        in place so custom metrics are returned and saved. \\
        Eight-generation training OOM & Used sequential microbatches of one
        while retaining all eight generations in each GRPO group. \\
        \bottomrule
    \end{tabular}
\end{table}

\section{Next Steps}

The immediate next step is to complete the full LA-CDM LoRA training run and
evaluate the final adapter on the unchanged 240-case test set. The trained
model will be compared with the Llama zero-shot baseline using diagnostic
accuracy, protocol-completion rate, macro F1, calibration error, test count,
and diagnostic cost.

Two ablation experiments are also planned: a Decision-Agent-only configuration
and an LA-CDM configuration without the test-cost reward. These comparisons
will help determine whether the Hypothesis Agent and cost-sensitive reward
improve accuracy, protocol adherence, calibration, or resource use.

\end{document}
